\documentclass[addpoints]{exam}
% \documentclass[addpoints,12pt]{exam}

\usepackage[slovene]{babel}
\usepackage{amsfonts}
\usepackage{amssymb}
\usepackage{amsmath}
\usepackage{float}
\usepackage{graphicx}
\usepackage{enumitem}
\usepackage{color}
\usepackage[gen]{eurosym}
\usepackage{newunicodechar}
\newunicodechar{€}{\euro}


%Komentar naslednje vrstice glede na verzijo testa -- rešitve ali prostor za odgovore:
% \printanswers

% \linespread{2}


\pointsinrightmargin
\marginbonuspointname{}


\renewcommand{\solutiontitle}{\noindent\textbf{Rešitve in točkovanje:}\par\noindent}

\colorgrids
\definecolor{GridColor}{gray}{0.7}
\setlength{\gridsize}{7mm}

\extrawidth{5mm}
\setlength{\rightpointsmargin}{1.5cm}

\begin{document}

\runningfooter{}{}{}

\begin{center}
    \fbox{\fbox{\parbox{15cm}{\centering
    Peto preverjanje znanja \\ 1. letnik \\ 21. maj 2025 \\ (Koordinatni sistem v ravnini, Funkcija, Premica)}}}
\end{center}

\vspace{0.1in}
\makebox[\textwidth]{Ime in priimek, oddelek:\enspace\hrulefill}


\begin{center}
    \hqword{Naloga}
    \hpword{Možne točke}
    \chbpword{Dodatne točke}
    \hsword{Dosežene točke}
    \htword{Skupaj}  
    \pointtable[h][questions]
    % \combinedgradetable[h][questions]

\end{center}

\vspace{0.1in}


\begin{questions}

    % 1. vprašanje

    \question
    Narišite dano množico točk v ravnini.
    \begin{parts}
        \part[4]
        $(3,8)\times (-\infty,2]$


        \part[4]
        $\left\{(x,y)\in\mathbb{R}^2 \mid (y=-2x+1)\land (y>-3)\right\} $

        \part[4]
        $(3,6)\times\{3,4,5\}$

    \end{parts}
    
    

    % 2. vprašanje
    \question[6]
    Izračunajte ploščino in obseg ter določite orientacijo trikotnika $\triangle ABC$ z oglišči v točkah $A(1,1)$, $B(3,5)$ in $C(-1,3)$.
    
        


    % 3. vprašanje
    \question
    Dana je premica z enačbo $y=\frac{2}{5}x+3$, ki seka koordinatni osi v točkah $M$ in $N$.

    \begin{parts}
        \part[4]
        Narišite premico.

            % \begin{solution}[6cm]
            %     ???
            % \end{solution}

        \part[6]
        Določite implicitno obliko enačbe vzporednice k dani premici, ki poteka skozi točko $(5,3)$.  Premico tudi narišite.

            % \begin{solution}[5cm]
            %     ???
            % \end{solution}

        \part[6]
        Določite segmentno obliko enačbe pravokotnice k dani premici, ki poteka skozi točko $(3,5)$.  Premico tudi narišite.

        \part[4]
        Naj bo $S$ razpolovišče daljice $MN$. Izračunajte razdaljo med točko $S$ in točko $(-2,1)$.

            % \begin{solution}[2cm]
            %     ???
            % \end{solution}

    \end{parts}
    


    % 5. vprašanje
    \question
    Dana je funkcija $$f(x)=\begin{cases}
    3x+2, & x\leq -1 \\ 2, & 0<x<3 \\ x-4, & x\geq 4
    \end{cases}.$$ 
    
    \begin{parts}
        \part[5]
        Narišite dano funkcijo.

            % \begin{solution}[9cm]
            %     ???
            % \end{solution}

        \part[3]
        Določite zalogo vrednosti in izračunajte ničli funkcije.

            % \begin{solution}[7cm]
            %     ???
            % \end{solution}

        \part[2]
        Določite in utemeljite, ali je funkcija surjektivna. 
        
        %     \begin{solution}[3cm]
        %         ???
        %     \end{solution}

        \part[1]
        Določite območje, kjer je funkcija naraščajoča.
        
            % \begin{solution}[2cm]
            %     ???
            % \end{solution}

    \end{parts}
    
    




    % 4. vprašanje
    \question
        Določite vrednost parametra $a$, da bo premica $(a-1)x+3y+6=0$:
        \begin{parts}
            \part[2] 
            potekala skozi točko $(2,6)$,

                % \begin{solution}[6cm]
                %     ???
                % \end{solution}


            \part[3]
            padajoča,
            
                % \begin{solution}[5cm]
                %     ???
                % \end{solution}

            \part[3] 
            imela diferenčni količnik enak $\frac{5}{3}$,
                        
                % \begin{solution}[6cm]
                %     ???
                % \end{solution}

            \part[3]
            vzporedna premici z enačbo $y=3x-2$.

            % \bonuspart[4]
            % \textit{Dodatna naloga:} oklepala s koordinatnima osema trikotnik s ploščino $6$.

        \end{parts}
    





\end{questions}



\end{document}